\documentclass[10pt,a4paper]{amsart}
\usepackage[margin=19mm]{geometry}
\usepackage{amsmath,amssymb,mathtools}
\usepackage[scr=rsfs,cal=euler]{mathalfa}
\usepackage{xfrac}
\usepackage[unicode,hidelinks,bookmarks=false]{hyperref}
\newcommand{\ie}{i.e.\ }
\newcommand{\cf}{cf.\ }
\newcommand{\resp}{respectively\ }
\newcommand{\Subsection}{\subsection}
\providecommand{\nobreakdash}{\nobreak\hbox{-}}
\setlength{\emergencystretch}{2em}
\multlinegap0pt
\usepackage{bm}
\usepackage{bbm}
\usepackage{dsfont}
\usepackage{xspace}
\usepackage{cancel}
\usepackage{mathtools}
\usepackage{amsthm}
\makeatletter
\@ifundefined{fact}{\newtheorem{fact}{Fact}}{}
\@ifundefined{definition}{\newtheorem{definition}{Definition}}{}
\@ifundefined{proposition}{\newtheorem{proposition}[fact]{Proposition}}{}
\makeatother
\makeatletter
\newcommand{\C}{\mathbb{C}}
\newcommand{\CC}{\ensuremath{\mathbb{C}}}
\newcommand{\Co}{\ensuremath{\mathbb{C}}}
\expandafter\ifx\csname claim\endcsname\relax
\newenvironment{claim}{\begin{cla} \it}{\end{cla}}
\fi
\expandafter\ifx\csname conjecture\endcsname\relax
\newenvironment{conjecture}{\begin{conj} \rm}{\end{conj}}
\fi
\expandafter\ifx\csname definition\endcsname\relax
\newenvironment{definition}{\begin{defi} \rm}{\end{defi}}
\fi
\newcommand{\eb}{\ensuremath{\mathbf{e}}}
\expandafter\ifx\csname example\endcsname\relax
\newenvironment{example}{\begin{exa} \rm}{\end{exa}}
\fi
\newcommand{\fb}{\ensuremath{\mathbf{f}}}
\newcommand{\gb}{\ensuremath{\mathbf{g}}}
\expandafter\ifx\csname observation\endcsname\relax
\newenvironment{observation}
%{\par\pushQED{\qed}\renewcommand\qedsymbol{$\diamond$}\nobs}
%{\popQED\endnobs}
\fi
\expandafter\ifx\csname remark\endcsname\relax
\newenvironment{remark}{\begin{rremark} \rm}{\end{rremark}}
\fi
\makeatother
\title[Triplets of Mutually Unbiased Bases]{Triplets of Mutually Unbiased Bases}
\author{M\'ate Matolcsi, \'Akos K. Matszangosz, D\'aniel Varga, Mih\'aly Weiner}
\date{Source: arXiv:2503.14752v2 (CC BY 4.0)}
\begin{document}
\maketitle

\noindent\emph{Source and licence.} Triplets of Mutually Unbiased Bases. Version arXiv:2503.14752v2 (CC BY 4.0), distributed under the Creative Commons Attribution 4.0 International licence, as recorded in the arXiv licence metadata for this exact version and shown on the abstract page https://arxiv.org/abs/2503.14752v2. Full rights notice: licenses/arxiv-2503.14752-CC-BY-4.0.txt.\par\smallskip

\noindent\textit{Editorial scope.} Counted unit: the Proposition whose source label is \texttt{misi} in the article (source0.tex lines 694--706, source label \texttt{misi}), delivered together with the article's own complete original proof of it (source0.tex lines 708--903, one \texttt{proof} environment of 196 lines). No mathematics is written, repaired or re-derived: statement and proof are the article's own text line for line. Required context retained uncounted: du (lines 373--377); hcube1 (lines 517--522); hcube (lines 544--555); prop:triplet\_equiv (lines 582--599); pq (lines 630--634); black\_and\_white (lines 688--691). Every definition, display and label of those lines is retained unchanged. Omitted from this package and not redistributed: 1--372, 378--516, 523--543, 556--581, 600--629, 635--687, 692--693, 707--707, 904--1667. Nothing inside a retained excerpt is omitted or altered: every retained body is delivered exactly as the article's lines are written, so no omission receipt of any kind accompanies this package. Citations: the retained excerpts carry no citation command at all, so no bibliography entry of the article is transcribed and no reference is redistributed. Reference adaptation: none was needed and none was made; every reference inside the delivered unit resolves inside it, so no unresolved reference or citation is produced. Printed slips kept literal: no printed word, symbol, hypothesis or number of the counted statement or of its proof was altered. Layout and class: the article's own class is \texttt{amsart} with packages including tikz, xcolor, inputenc, xy, lscape, enumerate, amsmath, amssymb, amsthm, overpic, hyperref, youngtab, cleveref, tikz-3dplot; the wrapper is the lane's pinned \texttt{amsart} editorial wrapper at 10pt on A4, which restates the theorem-like environments the retained text needs and adds the article's own macro definitions that the retained text needs (\texttt{\textbackslash C}, \texttt{\textbackslash CC}, \texttt{\textbackslash Co}, \texttt{\textbackslash eb}, \texttt{\textbackslash fb}, \texttt{\textbackslash gb}), with extra wrapper packages (bm, bbm, dsfont, xspace, cancel, mathtools). The delivered numbering is the unit's own. Assets: no retained span contains a figure environment, a graphics-inclusion command, a PDF-page inclusion, a table or a TikZ picture; no image file is copied into this package, the package declares no asset dependency and no image file is redistributed with it. Licence: version arXiv:2503.14752v2 of this article is distributed under the Creative Commons Attribution 4.0 International licence, as recorded in the arXiv licence metadata for this exact version and shown on the abstract page https://arxiv.org/abs/2503.14752v2. A full rights notice is delivered at licenses/arxiv-2503.14752-CC-BY-4.0.txt. Characters: every retained excerpt is pure ASCII, so the delivered engine, XeLaTeX with its default Latin Modern text font, reports no missing character.
\begin{definition}\label{du}
Suppose $X_1=(\eb^{(1)}_1,\ldots,\eb^{(1)}_d), \dots, X_m=(\eb^{(m)}_1,\ldots,\eb^{(m)}_d)$
 and $Y_1=(\fb^{(1)}_1,\ldots,\fb^{(1)}_d)$, $\dots$, $Y_m=(\fb^{(m)}_1,\ldots,\fb^{(m)}_d)$ are two systems of mutually unbiased bases. Let $P^{(k)}_j=|\eb^{(k)}_j\rangle \langle \eb^{(k)}_j|$ be the orthogonal projection onto the 1-dimensional subspace spanned by $\eb^{(k)}_j$ and, similarly, $\tilde{P}^{(k)}_j=|\fb^{(k)}_j\rangle \langle \fb^{(k)}_j|$.  
 We say that the two systems of MUBs $(X_1, \dots, X_m)$ and $(Y_1, \dots, Y_m)$ are {\it directly unitary equivalent} if there exists a unitary operator $U$ on $\Co^d$ such that $UP^{(k)}_jU^\ast = \tilde{P}^{(k)}_j$ for all $k \in [m]$ and $ j\in [d]$. 
\end{definition}

\begin{equation}\label{hcube1}
C_{j,k,l}= \sqrt{d^3}\,
\langle \eb_j, \fb_k\rangle 
\langle \fb_k, \gb_l\rangle
\langle \gb_l, \eb_j\rangle \;\;\;\; (1\le j,k,l\le d).
\end{equation}

\begin{proposition}\label{hcube}
For the Hadamard cube $C$ associated with a MUB-triplet $(X, Y, Z)$, as defined in \eqref{hcube1}, we have that
\begin{itemize}
    \item[i)] its entries are normalized: $\forall j,k,l\in[d]$: $|C_{j,k,l}|=1$,
    \item[ii)] each of its two-dimensional cross-sections ("slices") form a complex Hadamard matrix, i.e. for any fixed $j\in [d]$ the slices $C_{j,\cdot,\cdot}, C_{\cdot,j,\cdot}$ and $C_{\cdot,\cdot,j}$ are complex Hadamard matrices, 
    \item[iii)] parallel slices are phase-equivalent complex 
    Hadamard matrices; that is, for the parallel cross-sections $C_{j,\cdot,\cdot}$ and 
    $C_{j',\cdot,\cdot}$,  there exist scalars $\lambda_k, \mu_l$ ($k, l\in [d]$) such that $C_{j',k,l}=\lambda_k \, C_{j,k,l}\,\mu_l$, and similarly for parallel cross-sections in the other two directions,

    \item[iv)] its one-dimensional cross-sections ("piercings") sum to $\sqrt{d}$, that is, for any fixed $k,l\in[d]$ we have $\sum_{j=1}^d C_{j,k,l}=\sqrt{d}$ , and similarly for piercings of the other two directions.
\end{itemize}
\end{proposition}

\begin{proposition}
\label{prop:triplet_equiv}
Assume
$((\eb_j)_{j=1}^d,(\fb_j)_{j=1}^d,(\gb_j)_{j=1}^d)$ and
$((\tilde{\eb}_j)_{j=1}^d,(\tilde{\fb}_j)_{j=1}^d,(\tilde{\gb}_j)_{j=1}^d)$
are two MUB-triplets in $\C^d$. They
give rise to the same Hadamard cube, i.e. 
$$
\langle \eb_j, \fb_k\rangle 
\langle \fb_k, \gb_l\rangle
\langle \gb_l, \eb_j\rangle 
= 
\langle \tilde{\eb}_j, \tilde{\fb}_k\rangle 
\langle \tilde{\fb}_k, \tilde{\gb}_l\rangle
\langle \tilde{\gb}_l, \tilde{\eb}_j\rangle 
$$
for all $j,k,l\in [n]$, if and only if the MUB-triplets  are directly unitary equivalent in the sense of Definition \ref{du}.
\end{proposition}

\begin{equation}\label{pq}
A_{j,k} := | \eb_j\rangle\langle\eb_j|
\fb_k\rangle\langle\fb_k|=P_jQ_k\;\;\;\;
(j,k\in [d]).    
\end{equation}

\begin{equation}
\label{black_and_white}
C_{j',k,l} C_{j,k',l} C_{j,k,l'} C_{j',k',l'} = C_{j,k,l} C_{j',k',l} C_{j,k',l'} C_{j',k,l'}.
\end{equation}

\begin{proposition}\label{misi}
Assume $C=(C_{j,k,l})_{j,kl,\in [d]}$ is an array of complex numbers satisfying the following conditions:
\begin{itemize}
 \item[i')] its entries are normalized: $\forall j,k,l\in[d]$: $|C_{j,k,l}|=1$,
\item[ii')] at least one of its faces -- which we will call the ``bottom face'' -- is a complex Hadamard matrix, and the $d$ slices parallel to the bottom face are pairwise orthogonal, as elements of $\mathbb C^{d^2}$, 
\item[iii')] for any triplets of
indices $j,k,l$ and $j',k',l'$ equation (\ref{black_and_white}) holds,
\item[iv')]
entries in any one-dimensional ``horizontal'' cross-section (i.e.\ one parallel to the bottom face) 
sum to $\sqrt{d}$.
\end{itemize}
Then $C$ is a Hadamard cube, i.e. it satisfies properties $(i-iv)$ of Proposition \ref{hcube}. Also, there exists an MUB-triplet $(X, Y, Z)$ such that $C$ is the Hadamard cube associated with $(X, Y, Z)$.
\end{proposition} 

\begin{proof}
It is enough to show the last statement, i.e. that there exists an MUB-triplet associated to the cube $C$. Properties (i)-(iv) of Proposition \ref{hcube} will then follow automatically. 

\medskip

We can pick an arbitrary orthonormal basis $(\eb_j)_{j\in [d]}$ to be the first element of our MUB-triplet. To fix notations, let us say that the ``bottom face'' of $C$ is the one corresponding to the indeces $\{(j,k,l)\in [d]^3\,|\, l=d\}$. We then use this face to  define the second term of our MUB triplet; let
$\fb_j:= \frac{1}{\sqrt{d}} \sum_{k=1}^d C_{k,j,d}\eb_k$. The fact that $C_{\cdot,\cdot,d}$ is a complex Hadamard matrix implies that 
$(\fb_j)_{j\in [d]}$ is an orthonormal basis. Moreover, we have that $\langle \eb_k,\fb_j\rangle =\frac{1}{\sqrt{d}} C_{k,j,d}$; in particular, $\left((\eb_j)_{j\in [d]},(\fb_j)_{j\in [d]}\right)$
is an MUB pair. The  problem is to find a suitable third basis to form an MUB-triplet.  

\medskip

Let us look at this problem not in term of basis {\it vectors}, but -- as we have done before -- in terms of rank one orthogonal projections. Let $P_k:=|\eb_k\rangle\langle\eb_k|, Q_k:= |\fb_k\rangle\langle\fb_k| $ ($k\in [d]$) and $A_{j,k}:=P_jQ_k$ $(j,k\in [d])$ as in the proof of Prop.\ \ref{prop:triplet_equiv}. Instead of the vectors of the third basis, -- whose ``phases'', in any case, are not determined by our cube -- we need to find $d$ orthogonal projections $R_k$ $(k\in [d])$ such that
\begin{itemize}
\item[(a)] ${\rm Tr}(R_k R_j)=\delta_{k,j}$; this ensures that they are all rank one and are pairwise orthogonal -- i.e.\ they correspond to an ONB,
\item[(b)] ${\rm Tr}(P_k R_j)={\rm Tr}(Q_k R_j)=1/d$ for all $k,j\in [d]$; to ensure that the three bases form a MUB-triplet, and
\item[(c)] ${\rm Tr}(P_k Q_jR_l)= 1/\sqrt{d^3} \, C_{j,k,l}$ for all $j,k,l\in [d]$; to ensure that $C$ is indeed the Hadamard-cube corresponding to the constructed MUB-triplet.
\end{itemize}

As was already noted before, the matrices $A_{j,k}=P_jQ_k$ $(j,k\in [d])$ defined in \eqref{pq} form an orthogonal basis of $M_d(\CC)$, normalized to $1/\sqrt{d}$. It follows that with the definition
$$
R_l:=\frac{1}{\sqrt{d}}\sum_{j,k \in [d]}C_{j,k,l}\, A_{j,k}^* ={\sqrt{d}} \sum_{j,k\in [d]}C_{j,k,l}\, Q_kP_j\;\;\;
(l\in[d])
$$
we have that
$${\rm Tr}(P_j Q_kR_l)= 
{\rm Tr}(A_{j,k}R_l)=
\langle A_{j,k}^*,R_l\rangle_{H} = \frac{1}{\sqrt{d^3}} \, C_{j,k,l},$$
as required by point $(c)$. Moreover, we have 
\begin{eqnarray}
\nonumber
{\rm Tr}(Q_k R_l) &=& \frac{1}{\sqrt{d}}
\sum_{j',k'\in [d]}
 C_{j',k',l}
{\rm Tr}(Q_k A^*_{j',k'})=
\frac{1}{\sqrt{d}}
\sum_{j',k'\in [d]}
 C_{j',k',l}
{\rm Tr}(Q_k Q_{k'}P_{j'})
\\
\nonumber
&=&
\frac{1}{\sqrt{d}}
\sum_{j'\in [d]}
C_{j',k,l}
{\rm Tr}(Q_k P_{j'}) =
\frac{1}{d\sqrt{d}}
\sum_{j'\in [d]}
C_{j',k,l} =\frac{1}{d},
\end{eqnarray}
where the last equality follows by assumption (iv'). Similar computation shows that 
${\rm Tr}(P_k R_l)$ is also equal to $1/d$. Therefore,  point $(b)$ above is also satisfied.

\medskip

Since we have the coefficients of $R_l$ $R_{l'}$ in an orthogonal basis of $M_d(\CC)$, it is easy 
to compute the Hilbert-Schmidt product:
$$
\langle R_l,R_{l'}\rangle_{H}=
\frac{1}{d^2}\sum_{j,k} 
\overline{C}_{j,k,l} C_{j,k,l'} =
\delta_{l,l'},
$$
where we have used assumptions (i') and (ii'). Thus, ${\rm Tr}(R_l^*R_{l'})=\delta_{l,l'}$; which is almost what point $(a)$
asks for. What remains to be shown is that each $R_l$ is in fact an orthogonal projection; i.e.\ that
$R_l^* = R_l=R_l^2$. In order to proceed, let us note that
\begin{eqnarray}
\nonumber
A_{j,k}A_{j',k'} &=& 
|\eb_j\rangle\langle \eb_j|
\fb_k\rangle\langle \fb_k|
\eb_{j'}\rangle\langle \eb_{j'}|
\fb_{k'}\rangle\langle \fb_{k'}| 
=
\frac{\langle \eb_j,
\fb_k\rangle\langle \fb_k,
\eb_{j'}\rangle\langle \eb_{j'},
\fb_{k'}\rangle}{\langle \eb_j,\fb_{k'}\rangle}
\,
|\eb_j\rangle\langle \eb_{j}
|\fb_{k'}\rangle\langle \fb_{k'}|
\\
\nonumber
&=&
\frac{1}{d}\,
\frac{C_{j,k,d} \overline{C_{j',k,d}} C_{j',k',d}}{ C_{j,k',d}}
\,
A_{j,k'}
=
\frac{1}{d}\,
\frac{C_{j,k,d} C_{j',k',d}}{ 
C_{j',k,d}
C_{j,k',d}}
\,
A_{j,k'},
\end{eqnarray}
where we have used that for unit complex numbers conjugation is equivalent to taking reciprocals.
In particular, 
\begin{equation}
\label{eq:TrAA}
{\rm Tr}(A_{j,k}A_{j',k'})=
\frac{1}{d}\,
\frac{C_{j,k,d} C_{j',k',d}}{ 
C_{j',k,d}
C_{j,k',d}}
{\rm Tr}(P_jQ_{k'}) =
\frac{1}{d^2}\,\frac{C_{j,k,d} C_{j',k',d}}{ 
C_{j',k,d}
C_{j,k',d}}.
\end{equation}

We can now show that $R_l$ is an orthogonal projection. We have that
\begin{eqnarray}
\nonumber
(R_l^*)^2 &=& \left(\frac{1}{\sqrt{d}}\sum_{j,k \in [d]}\overline{C_{j,k,l}}\, A_{j,k}\right)^2 = 
\frac{1}{d}
\sum_{j,k,j',k'\in [d]}
\overline{C_{j,k,l}}\,
\overline{C_{j',k',l}}
A_{j,k}A_{j',k'}
\\
\nonumber
&=&
\frac{1}{d^2}
\frac{C_{j,k,d} \, C_{j',k',d}}{ 
C_{j,k,l}
C_{j',k',l} C_{j,k',d} C_{j',k,d}}
A_{j,k'} = 
\frac{1}{d^2}
\frac{C_{j,k,d} \, C_{j',k',d}}{ 
C_{j',k,l}
C_{j,k',l} C_{j,k,d} C_{j',k',d}}
A_{j,k'}
\\
\nonumber
&=&
\frac{1}{d^2}
\sum_{j,k,j',k'\in [d]}
\overline{C_{j',k,l}}\,
\overline{C_{j,k',l}}
A_{j,k'} = 
\left(\frac{1}{\sqrt{d^3}}
\sum_{j',k\in [d]}
\overline{C_{j',k,l}} \right)
\left(\frac{1}{\sqrt{d}}
\sum_{j,k'\in [d]}
\overline{C_{j,k',l}}
A_{j,k'} \right)
\\
\nonumber
&=&
\left(\frac{1}{\sqrt{d^3}}
\sum_{k\in [d]}
\sqrt{d}
\right) R_l^*
=R_l^*,
\end{eqnarray}
where from the second line to the third we used assumption (iii') -- i.e. the relation (\ref{black_and_white}) -- to ``reshuffle'' indices, and assumption (iv') from the third line to the fourth. Thus, the only thing that remains to be proven is the self-adjointness
of $R_l$.

To show that $R_l$ and $R_l^*$
coincide, it is enough to check that their Hilbert-Schmidt product with each $A_{j,k}^*$ ($j,k\in [d]$) coincide since the latter operators form an orthogonal basis of $M_d(\CC)$. 
We already know that
$\langle A_{j,k}^*,R_l\rangle_{H} = {\rm Tr}(P_j Q_k R_l)=1/\sqrt{d^3} C_{j,k,l}$. Also,
\begin{eqnarray}
\nonumber
\langle A_{j,k}^*,R_l^*\rangle_{H} &=& \frac{1}{\sqrt{d}}
\sum_{j',k'\in [d]} \overline{C_{j',k',l}}\,
{\rm Tr}(A_{j,k} A_{j',k'})=
\frac{1}{d^2\sqrt{d}}
\sum_{j',k'\in [d]} 
\frac{C_{j,k,d}\, C_{j',k',d}}{ 
C_{j',k',l} C_{j',k,d} C_{j,k',d}}
\\
\nonumber
&=&
\frac{1}{d^2\sqrt{d}}
\sum_{j',k'\in [d]} 
C_{j,k,l}
\overline{C_{j',k,l}}\, \overline{C_{j,k',l}}
=\frac{1}{d^2\sqrt{d}}
C_{j,k,l}
\sum_{j'\in [d]} 
\overline{C_{j',k,l}}\, 
\sum_{k'\in [d]} 
\overline{C_{j,k',l}}
\\
\nonumber
&=&
\frac{1}{\sqrt{d^3}}
C_{j,k,l}
\end{eqnarray}
where again from the first line to the second we used relation
(\ref{black_and_white}), and from 
the second line to the third assumption (iv'). The proof is now complete. 
\end{proof}

\end{document}
